% GENERATED FILE. Edit canonical lessons, then run npm run generate:books.
% canonical-source-hash: fnv1a64:57d1080f05f122f3
% canonical-lessons: GE-C356-karriere, GE-C356-aufstieg, GE-C356-herausforderung, GE-C356-bewaeltigen, GE-C356-scheitern

\chapter{Karriere machen}
\label{ch:ge-karriere-machen}
\hlchaptermodality{356}

\begin{chapteropening}
\textbf{By the end of this chapter:} \emph{I can talk about a career and a promotion, a challenge I coped with and a plan that failed.}

Complete the last lesson of chapter 356: I can talk about a career and a promotion, a challenge I coped with and a plan that failed.
\end{chapteropening}

\section[die Karriere]{die Karriere — the career}

\label{lesson:GE-C356-karriere}



\begin{quote}
Before the new one: say the German for the cooperation, then the German for to organize.
\end{quote}

\subsection*{You'll want to know: die Karriere}
\textbf{die Karriere} — \textquotedblleft{}the career\textquotedblright{}.

\begin{quote}
Sie war erst Kellnerin, dann machte sie Karriere.
\end{quote}

\begin{quote}
\emph{First she was a waitress, then she got on in her career.}
\end{quote}

Plural: \emph{die Karrieren}. \emph{Karriere machen} — to get on in your career.

\begin{cousinweb}
\emph{Karriere} was borrowed from French \emph{carrière} \textquotedblleft{}a road, a racecourse\textquotedblright{}, which came through Old Occitan \emph{carreira} from Late Latin \emph{carrāria} \textquotedblleft{}a road for carts\textquotedblright{} (EtymDB), built on \emph{carrus} \textquotedblleft{}a two-wheeled wagon\textquotedblright{} (Lewis \& Short): the track your working life runs along.
\end{cousinweb}

\begin{grammarlens}[title={What you've built}]
One more word for telling, explaining and describing at B1.
\end{grammarlens}

\subsection*{Your turn}
\begin{itemize}
\raggedright
  \item \emph{Say it:} \emph{die Karriere}
  \item \emph{Say it:} \emph{die Karriere}, once more
  \item \emph{Say it:} a sentence of your own with \emph{die Karriere}
  \item \emph{Recall:} say the German for the cooperation, then the German for to organize, then say \emph{die Karriere} again
\end{itemize}

\begin{tcolorbox}[breakable,colback=teal!4,colframe=teal!35!black,title={Before you move on}]
What does die Karriere mean? (\textquotedblleft{}The career\textquotedblright{}.) Say it once more.
\end{tcolorbox}
\section[der Aufstieg]{der Aufstieg — the rise, the promotion; the climb}

\label{lesson:GE-C356-aufstieg}



\begin{quote}
Before the new one: say the German for to organize, then the German for the career.
\end{quote}

\subsection*{You'll want to know: der Aufstieg}
\textbf{der Aufstieg} — \textquotedblleft{}the rise, the promotion; the climb\textquotedblright{}.

\begin{quote}
Nach fünf Jahren kam der Aufstieg zur Chefin.
\end{quote}

\begin{quote}
\emph{After five years came her promotion to boss.}
\end{quote}

Plural: \emph{die Aufstiege}. Also a climb up a mountain, and in sport the move up to a higher league.

\begin{cousinweb}
\emph{Aufstieg} goes with \emph{aufsteigen} \textquotedblleft{}to climb up\textquotedblright{}, built on \emph{steigen} \textquotedblleft{}to climb\textquotedblright{}, inherited from Middle High German \emph{stîgen}, Old High German \emph{stīgan} and Proto-Germanic \emph{stīganą} (EtymDB).
\end{cousinweb}

\begin{grammarlens}[title={What you've built}]
One more word for telling, explaining and describing at B1.
\end{grammarlens}

\subsection*{Your turn}
\begin{itemize}
\raggedright
  \item \emph{Say it:} \emph{der Aufstieg}
  \item \emph{Say it:} \emph{der Aufstieg}, once more
  \item \emph{Say it:} a sentence of your own with \emph{der Aufstieg}
  \item \emph{Recall:} say the German for to organize, then the German for the career, then say \emph{der Aufstieg} again
\end{itemize}

\begin{tcolorbox}[breakable,colback=teal!4,colframe=teal!35!black,title={Before you move on}]
What does der Aufstieg mean? (\textquotedblleft{}The rise, the promotion; the climb\textquotedblright{}.) Say it once more.
\end{tcolorbox}
\section[die Herausforderung]{die Herausforderung — the challenge}

\label{lesson:GE-C356-herausforderung}



\begin{quote}
Before the new one: say the German for the career, then the German for the rise.
\end{quote}

\subsection*{You'll want to know: die Herausforderung}
\textbf{die Herausforderung} — \textquotedblleft{}the challenge\textquotedblright{}.

\begin{quote}
Die neue Stelle ist eine große Herausforderung für mich.
\end{quote}

\begin{quote}
\emph{The new job is a big challenge for me.}
\end{quote}

Plural: \emph{die Herausforderungen}. From \emph{herausfordern} — to challenge.

\begin{cousinweb}
\emph{Herausforderung} is the noun of \emph{herausfordern} \textquotedblleft{}to challenge\textquotedblright{}, usually said to have first meant \textquotedblleft{}to call out\textquotedblright{} to a fight, from \emph{heraus} \textquotedblleft{}out\textquotedblright{} and \emph{fordern} \textquotedblleft{}to demand\textquotedblright{}. EtymDB has no history for \emph{fordern}; it is usually traced to the family of \emph{vorder} \textquotedblleft{}front\textquotedblright{}.
\end{cousinweb}

\begin{grammarlens}[title={What you've built}]
One more word for telling, explaining and describing at B1.
\end{grammarlens}

\subsection*{Your turn}
\begin{itemize}
\raggedright
  \item \emph{Say it:} \emph{die Herausforderung}
  \item \emph{Say it:} \emph{die Herausforderung}, once more
  \item \emph{Say it:} a sentence of your own with \emph{die Herausforderung}
  \item \emph{Recall:} say the German for the career, then the German for the rise, then say \emph{die Herausforderung} again
\end{itemize}

\begin{tcolorbox}[breakable,colback=teal!4,colframe=teal!35!black,title={Before you move on}]
What does die Herausforderung mean? (\textquotedblleft{}The challenge\textquotedblright{}.) Say it once more.
\end{tcolorbox}
\section[bewältigen]{bewältigen — to cope with, to manage (a task)}

\label{lesson:GE-C356-bewaeltigen}



\begin{quote}
Before the new one: say the German for the rise, then the German for the challenge.
\end{quote}

\subsection*{You'll want to know: bewältigen}
\textbf{bewältigen} — \textquotedblleft{}to cope with, to manage (a task)\textquotedblright{}.

\begin{quote}
Allein kann ich die viele Arbeit nicht bewältigen.
\end{quote}

\begin{quote}
\emph{I cannot cope with all this work on my own.}
\end{quote}

Inseparable: \emph{sie bewältigt, sie hat bewältigt}.

\begin{cousinweb}
\emph{bewältigen} is built on the family of \emph{Gewalt} \textquotedblleft{}power\textquotedblright{} and \emph{walten} \textquotedblleft{}to rule\textquotedblright{}, inherited from Old High German \emph{waltan} and Proto-Germanic \emph{waldaną} \textquotedblleft{}to rule\textquotedblright{} (EtymDB): you get power over a task.
\end{cousinweb}

\begin{grammarlens}[title={What you've built}]
One more word for telling, explaining and describing at B1.
\end{grammarlens}

\subsection*{Your turn}
\begin{itemize}
\raggedright
  \item \emph{Say it:} \emph{bewältigen}
  \item \emph{Say it:} \emph{bewältigen}, once more
  \item \emph{Say it:} a sentence of your own with \emph{bewältigen}
  \item \emph{Recall:} say the German for the rise, then the German for the challenge, then say \emph{bewältigen} again
\end{itemize}

\begin{tcolorbox}[breakable,colback=teal!4,colframe=teal!35!black,title={Before you move on}]
What does bewältigen mean? (\textquotedblleft{}To cope with, to manage (a task)\textquotedblright{}.) Say it once more.
\end{tcolorbox}
\section[scheitern]{scheitern — to fail}

\label{lesson:GE-C356-scheitern}



\begin{quote}
Before the new one: say the German for the challenge, then the German for to cope with.
\end{quote}

\subsection*{You'll want to know: scheitern}
\textbf{scheitern} — \textquotedblleft{}to fail\textquotedblright{}.

\begin{quote}
Der Plan ist am Geld gescheitert.
\end{quote}

\begin{quote}
\emph{The plan failed for lack of money.}
\end{quote}

Perfekt with \emph{sein}: \emph{er ist gescheitert}. \emph{An etwas scheitern} — to fail because of something.

\begin{cousinweb}
\emph{scheitern} is linked in EtymDB to Middle High German \emph{schīt} \textquotedblleft{}a split log\textquotedblright{}, today \emph{Scheit}. It is usually explained as first said of a ship that broke up into pieces of wood.
\end{cousinweb}

\begin{grammarlens}[title={What you've built}]
One more word for telling, explaining and describing at B1.
\end{grammarlens}

\subsection*{Your turn}
\begin{itemize}
\raggedright
  \item \emph{Say it:} \emph{scheitern}
  \item \emph{Say it:} \emph{scheitern}, once more
  \item \emph{Say it:} a sentence of your own with \emph{scheitern}
  \item \emph{Recall:} say the German for the challenge, then the German for to cope with, then say \emph{scheitern} again
\end{itemize}

\begin{tcolorbox}[breakable,colback=teal!4,colframe=teal!35!black,title={Before you move on}]
What does scheitern mean? (\textquotedblleft{}To fail\textquotedblright{}.) Say it once more.
\end{tcolorbox}
